\documentclass[11pt]{article}
\usepackage[margin=0.9in]{geometry}
\usepackage{fontspec}
\usepackage{amsmath,amssymb,amsthm}
\usepackage{xurl}
\begin{document}\sloppy
\begin{center}
{\Large GROK H0529: cube 32 of \texttt{a40ff54e580f5d0a} closes at 210 of 210 UNSAT; DRAT208; 56 of 75 kept cubes closed}\\[4pt]
2026-10-06 05:29 America/Los\_Angeles \quad branch \texttt{main}, tip \texttt{20674455} \quad automation 41774ea2\\[2pt]
\textbf{P versus NP: OPEN.}
\end{center}

\section*{Peer delta since H0515 (\texttt{8855aad1})}
Checked 05:30\,PT. Three commits landed, none under \path{research/main/} or \path{research/agents/grok/}; root \texttt{main.tex} is not touched.
\begin{itemize}
\item \texttt{381e5e7b}, LEAD-0004, \path{research/2026-10-06_0516_LEAD-0004_local_address_exponents.md}. It records LOCAL-ADDRESS-EXPONENT1 as PROVED/EXACT accounting: depth-first virtual recursion has time exponent $C=\max\bigl(r,(\eta+\lambda)/(1-\theta)\bigr)$ and space exponent $\sigma=\max(\ell,\mu)$, so a candidate aiming at the Williams SAT time-space endpoint must reach $C+\sigma<2\cos(\pi/7)$ after all conversions. LOCAL-ADDRESS-FORCING1 stays CANDIDATE/OPEN. I checked only the exponent algebra (below), not the model conversion or the stack terms, and I did not re-read Williams.
\item \texttt{03ac4668} and \texttt{20674455}, TASK-008 TWOSATDEC19 (\path{research/explorations/monotone-slice/2026-10-06_0523_TWOSATDEC19.md}): at $q=24$ the seed draws 16 SAT and 0 UNSAT, so the depth-$r$ residual is undefined; SAT-LOCAL-ERROR1 stays OPEN. I did not rerun it.
\end{itemize}
I do not upgrade or downgrade any peer label.

\section*{1. Literature}
None beyond the peer notes. The next claim is about our own finite search.

\section*{2. Claims recorded}
\textbf{Cube 32, $(0,1,3,6)$, closed (COMPUTATIONAL, solver-trusted).} The splitter log holds 210 distinct UNSAT subcube indices for cube 32, all 210 with the same step-2 tuple as the matching cube-0 subcube, no SAT line and no \texttt{TIMEOUT} line. The last one in was subcube 204, tuple $(1,4,6,14)$, UNSAT in 329.5\,s under the 600\,s cap; it is the slowest of the cube, and mean time is 30.2\,s. (The same tuple was cube 22's only timeout at 600\,s, later UNSAT at 674.5\,s under a 3000\,s rerun, so $(1,4,6,14)$ is the hard subcube in both.) The splitter printed \texttt{CUBE-DONE 32} at 05:31:43, which the closure check ignores.

\textbf{Claim DRAT208 (COMPUTATIONAL).} 208 of cube 0's 210 subcubes carry drat-trim-verified refutations (new: subcube 54). Remaining 2: subcubes 157 and 159; drat-trim on 157 has run since 05:28\,PT under a 3000\,s cap.

\textbf{Status of first-step cubes (COMPUTATIONAL, solver-trusted).} The closure list is now 0, 1, 2, 4, 9, 14, 19, 21, 22, 24, 27, 29, 30, 32. That makes 56 of 75 kept first-step cubes closed, 19 open, 0 SAT (up from 55 and 20 at H0515). Cube 33 is not a kept cube; the splitter moved on to cube 34, $(0,1,3,14)$, which had 25 UNSAT subcubes and no SAT at 05:33\,PT.

\section*{3. Smallest honest test}
\begin{itemize}
\item \textbf{Closure check} (\path{/workspace/scratch_h0529_1006/closure_run.py}, unchanged copy of H0348's; output \path{closure_0533.out}, after subcube 204 finished). It ignores \texttt{CUBE-DONE} and needs 210 distinct UNSAT indices with matching tuples. No SAT line; the only timeout is cube 22 subcube 204, superseded by its rerun; closed set as above (14 entries); partial \{34: 25, no open timeout\}. An earlier run at 05:30 (\path{closure_0530.out}) showed cube 32 at 209 with only subcube 204 missing.
\item \textbf{Cube scan} (\path{cscan.py 32}, output \path{cscan_0533.out}). 210 distinct, 210 of 210 tuples match, 0 missing, max 329.5\,s at subcube 204.
\item \textbf{Recount} (\path{/workspace/scratch_h2333_1005/experiments/recount.py}, rerun 05:33\,PT). Total verified 208, remaining 2, 127 distinct CNF hashes among 127 batch lines, none not verified.
\item \textbf{LEAD-0004 exponent algebra} (\path{lead4_fp.py}, output \path{lead4_fp.out}). On 20000 random parameter draws ($r\in[1,3]$, $\eta,\lambda\in[0,1]$, $\theta\in[0,0.95]$), $C=\max(r,(\eta+\lambda)/(1-\theta))$ is a fixed point of $C=\max(r,\eta+\lambda+\theta C)$ and the iteration from $C=r$ converges to it: 0 mismatches. Recomputed $2\cos(\pi/7)=1.801938$. This checks only that the closed form solves the exponent recurrence $T(N)=N^r+N^{\eta+\lambda}T(N^\theta)$; that this recurrence is the right charge for the model is not checked here.
\end{itemize}

\section*{4. What this changes}
Closed rises from 55 to 56 of 75; certified rises from 207 to 208 of 210. Neither moves a charged resource. The bounds stay $9\le L(\texttt{a40ff54e580f5d0a},2)\le 11$ (COMPUTATIONAL for the upper bound; the lower bound stays CONDITIONAL on the tightness lemma as recorded earlier); $L=9$ itself stays open while 19 cubes are open. The largest confirmed price stays one.

\section*{Effect on the standing candidate}
No claim is falsified and no charged resource improves. The time-state-representation surface stays CANDIDATE. LEAD-0004 is accounting for virtual recursion and, like PROOF-0011, prices forget-then-recompute; it constrains a candidate but does not supply one, so it does not touch the open arrows. The SAT-to-charged-model compiler is OPEN, and so is the product-span obstruction. The finite results concern one fixed function and imply no lower bound for any family.

\section*{Next smallest step}
Let cube 34 finish its 210 subcubes and run the closure check; watch subcube 204 $(1,4,6,14)$ in particular and rerun any timeout at 3000\,s. Let the DRAT batch finish 157 and 159.

\bigskip
\noindent\textbf{Status sentence.} P versus NP remains OPEN: the SAT-to-charged-model compiler and the product-span obstruction are open load-bearing arrows. This pass records that cube 32 of \texttt{a40ff54e580f5d0a} at tight $L=9$ is closed with 210 of 210 subcubes UNSAT and no timeout (COMPUTATIONAL, solver-trusted), and that 208 of cube 0's 210 subcubes carry drat-trim-verified refutations (COMPUTATIONAL); 56 of 75 kept first-step cubes are closed, 19 open, 0 SAT, and $9\le L(f,2)\le 11$ is unchanged. No claim is falsified, no charged resource improves, and this is not a proof of anything about P versus NP.
\end{document}
